\scanpage{167}
% archive foliation: 166
\uline{Ex} Si $X$ \deleted{\unsure{est intègre les}} \inserted{est \uline{réduit}, alors schématiquement dense $=$ dense} \deleted{\unsure{ouverts sont les ouverts non vides}}

pseudo-morphismes $=$ applications rationnelles

Pour tout ouvert $U$ de $X$, soit $F(U)$ l'ens des pseudo-$S$-morphismes de $U$ dans $Y$ ($S$-schéma fixé). Applications de restriction $F(U)\unsure{-|\to}F(V)$ si $U\supset V$. Cela \deleted{est} définit un \uline{faisceau} [grâce à la relation $\operatorname{Dom}(f_i|U_{ij})=\operatorname{Dom}(f_j|U_{ij})$], appelé \uline{faisceau des} \deleted{\unsure{quasi-morphismes}} \uline{pseudo-$S$-morphismes} de $X$ ds $Y$.

En particulier, on obtient ainsi le faisceau des pseudo-sections de $O_X$ (faisant $X=S$, $Y=X[t]$) qui est un faisceau d'anneaux $\supset\underline O_X$, en particulier un faisceau d'algèbres.

\deleted{\uline{Prop.} 0.1 Ce faisceau est quasi-cohérent.
\unsure{Il suffit pour} $X=\operatorname{Spec}(A)$.}

